\section*{Mit diesem Übungsheft arbeiten}
Dieses Heft ergänzt die
\href{_B10-binomial-reading.pdf\#binomial.reading}{Lesefassung
\emph{Binomialkoeffizienten und das Pascalsche Dreieck}}.
Acht Aufgaben führen vom Zeichnen einzelner Wege zum Begründen der
Formeln. Die vollständigen Lösungen beginnen auf Seite~\pageref{bin:solutions}.

\textbf{Unser Wegemodell.} Wir starten an der Spitze $(0,0)$.
Jeder Schritt führt schräg nach links unten ($L$) oder rechts unten ($R$):
\[
 L:(n,k)\longmapsto(n+1,k),\qquad
 R:(n,k)\longmapsto(n+1,k+1).
\]
Nach $n$ Schritten, davon $k$ nach rechts, stehen wir bei $(n,k)$.
Dorthin führen $\binom nk$ Wege. Der obere Index zählt alle Schritte,
der untere die Rechtsschritte. Ein Weg ohne Schritte zählt als ein Weg.

\textbf{So arbeitest du.} Zeichne oder notiere zunächst Beispiele.
Begründe dann, warum deine Einteilung jeden Weg genau einmal erfasst.
Bei einem Produkt erkläre, welche Möglichkeiten unabhängig kombiniert
werden; bei einer Summe, welche getrennten Klassen zusammenkommen.
Nutze die Karoflächen rechts und unter den Aufgaben. Für die Beweise
in Aufgabe~6 und~8 gibt es jeweils eine zusätzliche Arbeitsseite mit Hinweis.

Wie in der Lesefassung gilt $\N=\{0,1,2,\ldots\}$,
$\binom nk=0$ für $k>n$ und $a^0=1$, auch für $a=0$.
Stellen innerhalb einer Zeile zählen wir von links ab $0$.

\medskip\tableofcontents
